

\section{Le "Monitor" de Krashen : La grammaire consciente ne produit pas la fluidité}
space{0.3cm}
oindent

L'enseignement du grec ancien, discipline séculaire et pilier des humanités, traverse une crise
structurelle profonde concernant ses résultats pédagogiques. Malgré des années d'étude rigoureuse,
une majorité significative d'étudiants — et parfois même d'enseignants — peine à atteindre une
véritable compétence de lecture fluide, définie comme la capacité de comprendre un texte directement
dans la langue cible sans recourir à une traduction mentale intermédiaire ou à un décodage
analytique laborieux.\cite{II-2-3-ref1} Ce phénomène, souvent accepté avec fatalisme comme une
conséquence inhérente à la difficulté des langues classiques, mérite d'être réexaminé à la lumière
des théories contemporaines de l'acquisition des langues secondes (SLA) et des neurosciences
cognitives.

Ce rapport se concentre spécifiquement sur le point II.2.\cite{II-2-3-ref3} de votre plan d'étude :
\textit{\og Le 'Monitor' de Krashen : la grammaire consciente comme correcteur lent, insuffisante
pour la production ou la lecture fluide\fg{}}. L'objectif est de démontrer, par une analyse
exhaustive de la littérature scientifique, que le recours prédominant à la grammaire explicite (le
savoir conscient sur la langue) constitue non pas une aide, mais un goulot d'étranglement cognitif
qui empêche l'accès à la fluidité.

Nous nous appuierons sur l'hypothèse du Monitor formulée par Stephen Krashen, que nous croiserons
avec le modèle neurobiologique Déclaratif/Procédural de Michael Ullman et les recherches sur la
charge cognitive et la mémoire de travail. Cette triangulation théorique permettra de mettre en
évidence les mécanismes psycholinguistiques qui rendent la méthode traditionnelle de \og Grammaire-
Traduction\fg{} structurellement inapte à produire des lecteurs compétents en grec ancien, confinant
les apprenants dans un statut perpétuel de \og déchiffreurs\fg{}.\cite{II-2-3-ref2}



\subsection{L'architecture théorique : l'hypothèse du monitor et la dichotomie acquisition/apprentissage}

Pour comprendre pourquoi la grammaire consciente échoue à générer de la fluidité, il est impératif
de déconstruire l'architecture cognitive proposée par la théorie de l'acquisition des langues de
Krashen. Cette théorie ne rejette pas la grammaire \textit{per se}, mais elle lui assigne une
fonction cognitive très spécifique et limitée : celle de \og Monitor\fg{}.



\subsection{La distinction fondamentale : acquisition vs apprentissage}

La pierre angulaire de l'argumentation réside dans la distinction stricte entre deux processus
cognitifs indépendants d'internalisation d'une langue. Cette distinction, souvent ignorée dans la
pédagogie classique qui tend à les confondre, est pourtant cruciale pour comprendre les échecs de la
méthode traditionnelle.

L'\textbf{Acquisition} est un processus subconscient, intuitif et implicite. Il est analogue à la
manière dont les enfants développent leur langue maternelle. L'acquisition se produit lorsque
l'apprenant est focalisé sur le sens du message, sur la communication réelle, et non sur la forme
linguistique. C'est ce système acquis qui est responsable de la fluidité, de la spontanéité et de
l'intuition grammaticale (le \og sentiment\fg{} qu'une phrase sonne juste ou
faux).\cite{II-2-3-ref4}

L'\textbf{Apprentissage}, en revanche, est un processus conscient, explicite et formel. Il résulte
de l'enseignement direct, de l'étude des règles grammaticales, de la mémorisation des paradigmes de
déclinaisons et de conjugaisons. C'est le domaine du \og savoir sur la langue\fg{}
(métalinguistique). Selon Krashen, ce système appris ne se transforme jamais en système acquis
(position de non-interface). Par conséquent, les règles apprises consciemment ne peuvent pas générer
de la parole ou de la compréhension fluide directement.\cite{II-2-3-ref6}



\subsection{Le rôle restreint du monitor}

Dans ce cadre, le système appris n'a qu'une seule fonction possible : celle de \textbf{Monitor} (ou
d'éditeur). Le Monitor ne peut pas initier la production d'une phrase ni le décodage rapide du sens.
Son rôle est d'intervenir \textit{a posteriori} (ou juste avant l'émission) pour scanner la
production du système acquis, la comparer aux règles conscientes stockées en mémoire, et effectuer
des corrections si nécessaire.\cite{II-2-3-ref8}

Cette fonction de surveillance est intrinsèquement parasitaire sur le flux de la communication. Elle
demande à l'apprenant de suspendre temporairement son attention au sens pour la rediriger vers la
forme. Dans le contexte du grec ancien, cela se traduit par l'étudiant qui, au milieu d'une phrase
de Thucydide, cesse de suivre le raisonnement historique pour analyser si un verbe est un aoriste ou
un imparfait, brisant ainsi la continuité cognitive nécessaire à la compréhension profonde.



\subsection{Les trois conditions de l'activation du monitor}

Krashen a identifié trois conditions sine qua non pour que le Monitor puisse être utilisé
efficacement. L'analyse de ces conditions dans le contexte spécifique du grec ancien révèle pourquoi
cette stratégie est vouée à l'échec pour la lecture fluide.\cite{II-2-3-ref8}

\begin{table}[h!]\small\centering\begin{tabular}{| p{2.3cm} | p{3.5cm} | p{3.5cm} |}\hline\textbf{Condition} & \textbf{Description Théorique} & \textbf{Obstacle Majeur en Grec Ancien} \\ \hline\textbf{1. Temps (Time)} & L'utilisateur doit disposer d'un temps suffisant pour récupérer la règle et l'appliquer. & La lecture fluide ou l'écoute d'un texte exige un traitement en temps réel (millisecondes). S'arrêter pour analyser chaque mot (morphologie flexionnelle) ralentit la vitesse de lecture en deçà du seuil de compréhension (voir section IV). \\ \hline\textbf{2. Focalisation sur la Forme (Focus on Form)} & L'attention doit être dirigée vers la correction grammaticale (\textit{correctness}) plutôt que sur le contenu. & Se concentrer sur les désinences (ex: génitif vs datif) détourne les ressources attentionnelles du sens narratif ou philosophique du texte, rendant la lecture \og intelligente\fg{} impossible. \\ \hline\textbf{3. Connaissance de la Règle (Know the Rule)} & L'utilisateur doit posséder une représentation mentale explicite et correcte de la règle applicable. & La complexité extrême du grec (dialectes, verbes irréguliers, règles d'accentuation, sandhi) fait que même les experts ne \og connaissent\fg{} pas consciemment toutes les règles à tout moment. La règle est souvent trop complexe pour une application consciente rapide. \\ \hline\end{tabular}\end{table}Ces conditions limitent drastiquement l'utilité du Monitor. Krashen suggère que le Monitor n'est
réellement utile que dans des situations d'écriture formelle où le temps n'est pas contraint, ou
lors d'examens de grammaire.\cite{II-2-3-ref11} Tenter de l'utiliser pour la lecture courante
revient à essayer de courir un marathon en analysant consciemment la biomécanique de chaque foulée :
c'est inefficace, épuisant et lent.



\subsection{Fondements neurocognitifs : le modèle déclaratif/procédural}

L'hypothèse du Monitor, formulée initialement sur la base d'observations comportementales, trouve
aujourd'hui une validation robuste et une explication mécaniste dans les neurosciences cognitives,
notamment à travers le modèle Déclaratif/Procédural (DP Model) développé par Michael Ullman. Ce
modèle fournit le substrat biologique expliquant pourquoi la grammaire consciente est un correcteur
lent.\cite{II-2-3-ref13}



\subsection{La dualité des systèmes de mémoire}

Le modèle DP postule que le langage repose sur deux systèmes de mémoire cérébraux distincts, qui
correspondent remarquablement à la distinction Acquisition/Apprentissage de
Krashen.\cite{II-2-3-ref15}

1. \textbf{La Mémoire Déclarative (Le Siège du Monitor) :}

\begin{itemize}\item \textbf{Localisation :} Lobe temporal médian (incluant l'hippocampe) et régions néocorticales associées.\item \textbf{Fonction :} Stockage des faits, des événements, et des connaissances explicites (\og Savoir que\fg{}). Dans le contexte linguistique, elle gère le lexique (associations arbitraires forme-sens) et les règles grammaticales apprises explicitement.\item \textbf{Caractéristiques de Traitement :} Apprentissage rapide (parfois en une seule exposition), mais récupération consciente, attentionnelle et relativement lente. Elle est flexible mais coûteuse en ressources cognitives.\item \textbf{Application au Grec :} C'est le système utilisé lorsqu'un élève récite un tableau de déclinaison (\textit{logos, logou, logô...}) ou applique consciemment la règle de la contraction vocalique.\end{itemize}2. \textbf{La Mémoire Procédurale (Le Siège de l'Acquisition) :}

\begin{itemize}\item \textbf{Localisation :} Circuits fronto-striataux, incluant les ganglions de la base (noyau caudé, putamen) et l'aire de Broca (Brodmann 44/45).\item \textbf{Fonction :} Apprentissage et exécution d'habiletés motrices et cognitives (\og Savoir comment\fg{}). Elle gère la grammaire internalisée, la syntaxe automatique et les séquences régulières.\item \textbf{Caractéristiques de Traitement :} Apprentissage lent et graduel nécessitant une exposition répétée, mais exécution rapide, automatique et inconsciente. Elle fonctionne en parallèle et demande peu de ressources attentionnelles une fois la compétence acquise.\item \textbf{Application au Grec :} C'est le système qui permet de comprendre instantanément que \textit{élabon} est un aoriste sans avoir à décomposer le mot, simplement parce que le cerveau reconnaît le \og motif\fg{} neuronal.\end{itemize}

\subsection{La latence différentielle et les temps de réaction}

La raison pour laquelle le Monitor est \og insuffisant pour la fluidité\fg{} est physiologiquement
ancrée dans la différence de vitesse de traitement entre ces deux systèmes. Les études sur les temps
de réaction (Reaction Time \- RT) montrent systématiquement que l'accès aux connaissances explicites
(déclaratives) est significativement plus lent que le traitement implicite
(procédural).\cite{II-2-3-ref17}

Le traitement procédural est un réflexe cognitif : le stimulus (le mot grec) déclenche directement
la représentation grammaticale. Le traitement déclaratif, lui, nécessite une recherche active en
mémoire, souvent médiée par la mémoire de travail. Pour un apprenant de grec s'appuyant sur le
Monitor, chaque mot déclenche une routine de vérification consciente :

1. Perception visuelle du mot.

2. Isolation de la désinence.

3. Recherche dans la \og bibliothèque\fg{} déclarative de la règle correspondante.

4. Vérification de la compatibilité (genre, nombre, cas).

5. Attribution du sens.

Ce processus sériel prend plusieurs centaines de millisecondes, voire plusieurs secondes par mot. À
l'échelle d'une phrase complexe, ces délais s'accumulent, créant une latence incompatible avec le
maintien du sens global en mémoire de travail.\cite{II-2-3-ref19}



\subsection{L'effet de balancier ("see-saw effect") : compétition neuronale}

Un aspect critique du modèle d'Ullman, qui a des implications pédagogiques majeures pour le grec
ancien, est l'effet de \og balancier\fg{} (\textit{See-Saw Effect}). Les recherches suggèrent qu'il
existe une relation compétitive, voire inhibitrice, entre la mémoire déclarative et la mémoire
procédurale. L'activation intense de l'une peut supprimer ou atténuer l'activité de
l'autre.\cite{II-2-3-ref16}

Cela signifie que l'enseignement traditionnel (Grammaire-Traduction), qui force l'élève à utiliser
intensivement sa mémoire déclarative (analyser, décortiquer, justifier chaque forme grammaticale),
pourrait physiologiquement empêcher le système procédural de prendre le relais. En surentraînant le
Monitor, on consolide les circuits de l'hippocampe au détriment des circuits striataux nécessaires à
la fluidité. L'élève devient un expert en règles (déclaratif) mais reste un lecteur inefficace
(procédural déficient).\cite{II-2-3-ref23}



\subsection{La surcharge cognitive : pourquoi le grec ancien résiste au monitor}

L'inefficacité du Monitor est exacerbée par la nature spécifique de la langue grecque. Si le Monitor
peut éventuellement gérer des langues à morphologie simple (comme l'anglais ou l'espagnol de base)
avec un ralentissement modéré, il s'effondre face à la complexité flexionnelle du grec ancien.



\subsection{La complexité morphologique et la règle du "know the rule"}

La condition \og Know the Rule\fg{} de Krashen est particulièrement problématique ici. La \og
règle\fg{} en grec ancien n'est pas une simple relation binaire. Elle implique souvent une
combinatoire complexe de facteurs phonétiques, morphologiques et syntaxiques.

Par exemple, pour analyser correctement une forme verbale, l'élève doit gérer :

\begin{itemize}\item L'augment (temporel ou syllabique).\item Le redoublement (avec ses variations selon la consonne initiale).\item Les modifications vocaliques (alternances, contractions).\item Les désinences primaires vs secondaires.\item Les variations dialectales (ionien vs attique dans les textes littéraires).\end{itemize}Traiter cette masse d'informations via le Monitor sature instantanément la \textbf{mémoire de
travail}. La mémoire de travail humaine a une capacité limitée (environ 4 à 7 éléments). Si
l'analyse d'un seul verbe occupe 3 ou 4 \og slots\fg{} de mémoire de travail pour gérer les règles
morphologiques, il ne reste plus d'espace pour retenir le sujet de la phrase, le contexte narratif
ou les nuances sémantiques.\cite{II-2-3-ref19} C'est la surcharge cognitive : le processeur central
est bloqué à 100\% sur le décodage local, rendant l'intégration globale impossible.



\subsection{L'illusion de la maîtrise par la gtm}

La méthode Grammaire-Traduction (GTM) repose sur l'idée que la maîtrise consciente de ces règles
mène à la compétence. Cependant, les recherches montrent que cette approche produit typiquement des
\textbf{\og Monitor Over-users\fg{}} (Sur-utilisateurs du Monitor).\cite{II-2-3-ref26}

Ces apprenants se caractérisent par :

\begin{itemize}\item Une hésitation constante à l'oral et une lecture hachée.\item Une anxiété liée à la correction grammaticale (Filtre Affectif élevé).\item Une incapacité à traiter le texte comme un flux de sens.\item Une dépendance à la traduction mot-à-mot.\end{itemize}Le sur-utilisateur du Monitor en grec ancien aborde le texte comme un champ de mines grammaticales.
Il ne lit pas pour savoir ce qu'Achille a dit à Agamemnon ; il lit pour identifier un accusatif ou
un optatif. Cette approche analytique crée une \og compétence apparente\fg{} (l'élève réussit les
tests de grammaire) mais masque une incompétence fonctionnelle (l'élève ne peut pas lire couramment
sans dictionnaire ni grammaire de référence).\cite{II-2-3-ref28}



\subsection{Interférence et charge de la mémoire de travail}

Les études utilisant des potentiels évoqués (ERP) montrent que les violations syntaxiques traitées
par des natifs (traitement procédural) déclenchent une signature cérébrale spécifique (LAN/P600)
très rapide. Chez les apprenants tardifs s'appuyant sur des règles explicites (traitement
déclaratif), ces signatures sont souvent absentes, retardées, ou remplacées par des activations
massives indiquant un effort cognitif intense et conscient.\cite{II-2-3-ref29}

En grec ancien, où l'ordre des mots est libre et où les accords (adjectif-nom) peuvent être séparés
par plusieurs propositions (hyperbates), la charge sur la mémoire de travail est immense. Le lecteur
doit \og garder en attente\fg{} un adjectif au nominatif pendant qu'il traite trois lignes de texte
avant de trouver son substantif. Le système procédural (acquis) gère cela efficacement par des
attentes probabilistes inconscientes. Le Monitor (appris), en revanche, doit maintenir activement
cette information en mémoire consciente, ce qui est extrêmement coûteux et fragile. Dès que
l'attention faiblit ou qu'un mot inconnu surgit, la chaîne de retenue se brise et le sens
s'effondre.\cite{II-2-3-ref31}



\subsection{Fluidité de lecture : seuils et limitations du décryptage}

La lecture fluide n'est pas simplement une version accélérée du décryptage grammatical. C'est un
processus qualitativement différent qui requiert des vitesses de traitement que le Monitor ne peut
physiquement pas atteindre.



\subsection{La vitesse de lecture comme indicateur de processus}

Les recherches de Paul Nation et d'autres sur la lecture en langue seconde établissent qu'une
lecture confortable et compréhensive se situe généralement autour de \textbf{200 à 250 mots par
minute (wpm)}. En dessous d'un certain seuil (environ 100-150 wpm), la mémoire tampon phonologique
s'efface avant que le sens de la phrase ne soit construit. Le lecteur oublie le début de la phrase
avant d'en atteindre la fin.\cite{II-2-3-ref32}

Le traitement conscient via le Monitor est intrinsèquement lent. Un \og déchiffreur\fg{} de grec
utilisant la méthode analytique plafonne souvent à \textbf{10 ou 20 mots par minute}. À cette
vitesse, la lecture est cognitivement impossible : c'est une activité de résolution de puzzle, pas
de réception de message. Le Monitor agit comme un frein permanent, empêchant l'automatisation
nécessaire pour atteindre les vitesses où la compréhension globale émerge.\cite{II-2-3-ref34}



\subsection{Le seuil lexical de 98\% et l'impuissance du monitor}

Une autre donnée critique fournie par les recherches de Nation et Laufer est le seuil de couverture
lexicale. Pour qu'une lecture soit fluide et permette l'apprentissage incident (deviner le sens par
le contexte), le lecteur doit connaître \textbf{95\% à 98\%} des mots du texte.\cite{II-2-3-ref35}

\begin{table}[h!]\small\centering\begin{tabular}{| p{2.3cm} | p{3.5cm} | p{3.5cm} |}\hline\textbf{Couverture Lexicale} & \textbf{Impact sur la Compréhension} & \textbf{Rôle du Monitor} \\ \hline\textbf{98-100\%} & Lecture fluide, plaisir, apprentissage incident possible. & Inactif ou utilisé très rarement pour des nuances stylistiques. \\ \hline\textbf{95-98\%} & Compréhension globale bonne, quelques arrêts. & Activé ponctuellement pour vérifier une hypothèse contextuelle. \\ \hline\textbf{\< 90\%} & \og Frustration level\fg{}, décryptage laborieux, perte du fil narratif. & \textbf{Sur-activé.} Le Monitor tente de compenser les lacunes lexicales par l'analyse grammaticale, menant à la surcharge cognitive. \\ \hline\end{tabular}\end{table}Dans l'enseignement traditionnel du grec, les élèves sont souvent confrontés dès le début à des
textes authentiques où leur couverture lexicale est bien inférieure à 80\%. Ils sont forcés de
compenser ce déficit lexical par une analyse syntaxique hyper-intensive via le Monitor. Mais comme
le montrent Hu et Nation (2000), aucune quantité d'analyse grammaticale (Monitor) ne peut compenser
un déficit lexical massif pour la compréhension fluide. Le Monitor ne peut pas \og inventer\fg{} le
sens d'un mot inconnu ; il ne peut qu'analyser sa fonction.\cite{II-2-3-ref38}



\subsection{L'inhibition de l'intuition probabiliste}

La fluidité repose sur la prédiction. Le cerveau expert anticipe le mot suivant ou la structure
grammaticale à venir (traitement top-down). Cette anticipation est probabiliste et basée sur
l'expérience (input massif). Le Monitor, en revanche, fonctionne en mode bottom-up strict : il
attend la donnée, l'analyse, et valide.

L'utilisation excessive du Monitor inhibe le développement de ces mécanismes prédictifs. Au lieu de
laisser le cerveau deviner que tèn sera suivi d'un substantif féminin accusatif et de pré-activer
les candidats sémantiques probables, le Monitor Over-user attend passivement de voir le mot pour
vérifier son accord. Cette passivité prédictive est fatale pour la lecture fluide de textes longs
comme les épopées homériques ou les dialogues platoniciens.\cite{II-2-3-ref20}



\subsection{Implications pédagogiques : vers une approche post-monitor}

L'analyse convergente de l'hypothèse de Krashen, du modèle DP d'Ullman et des recherches
psycholinguistiques mène à une conclusion sans appel : la grammaire consciente, lorsqu'elle est
positionnée comme le mécanisme central de l'apprentissage (comme dans la GTM), est un obstacle à la
fluidité en grec ancien.



\subsection{Repenser le rôle de la grammaire}

L'objectif ne doit pas être d'éliminer toute grammaire, mais de la remettre à sa place physiologique
:

\begin{itemize}\item \textbf{La Grammaire comme Vérificateur Tardif :} Le Monitor est utile pour l'édition de production écrite (thème grec) ou pour l'analyse littéraire fine d'un texte \textit{déjà compris}. Il ne doit pas être l'outil de découverte du texte.\item \textbf{La Grammaire \og Pop-up\fg{} :} Au lieu de lourdes leçons explicites \textit{a priori}, des explications grammaticales brèves et contextuelles (moins de 10 secondes) peuvent satisfaire la curiosité sans enclencher le mode déclaratif lourd ni l'effet de balancier inhibiteur.\cite{II-2-3-ref12}\end{itemize}

\subsection{Priorité à l'input compréhensible}

Pour développer le système procédural (Acquisition) capable de traiter la complexité du grec en
temps réel, la seule voie validée par la recherche SLA est l'exposition massive à un \textbf{Input
Compréhensible} (\textit{Comprehensible Input}).

\begin{itemize}\item L'enseignement doit fournir des messages que l'élève comprend, focalisés sur le sens.\item La lecture doit être extensive et adaptée au niveau (lectures graduées), respectant le seuil de 98\% de vocabulaire connu, pour permettre l'automatisation des processus neuronaux.\item L'évaluation ne doit pas porter sur la capacité à réciter des règles (Monitor) mais sur la capacité à comprendre et interagir avec le texte (Acquisition).\cite{II-2-3-ref1}\end{itemize}

\subsection{Conclusion}

En somme, le Monitor de Krashen est un \og correcteur lent\fg{} car il utilise des circuits
neuronaux (mémoire déclarative) qui ne sont pas conçus pour la vitesse de traitement requise par le
langage naturel. Dans le contexte du grec ancien, l'obstination pédagogique à passer par ce canal
étroit et coûteux explique la prévalence de l'échec en lecture fluide. Pour former des lecteurs
compétents, la didactique du grec doit opérer une révolution copernicienne : cesser d'entraîner le
correcteur interne (le Monitor) et commencer à nourrir le système d'acquisition implicite par un
input riche, abondant et compréhensible. Tant que la grammaire consciente restera le roi, la
fluidité restera l'exilée.

